\chapter{Fine-Tuning for Classification}\label{ch:classify}

A pretrained language model is a general-purpose engine for predicting the next token.
Most deployed models, though, have a narrower job: route a support ticket, flag a
message as spam, decide whether a code change touches a security-sensitive file.
\emph{Classification fine-tuning} is the most direct way to turn a pretrained model into
such a component. We keep the transformer trunk, throw away the 50{,}257-way vocabulary
projection, put a tiny $C$-way projection in its place, and train on labelled
examples.

The modelling change is small, but it raises several systems questions that recur for
the rest of the book. Which parameters do we update, and what does freezing the rest
save in memory and in FLOPs? How much of each batch is padding, and how do we stop
paying for it? Which token's hidden state should the head read, and does that answer
change with padding? A classifier is also the simplest example of a \emph{prefill-only}
inference workload: one forward pass, no decoding loop, no KV cache that outlives the
request. That makes it a useful contrast to the generative serving systems of
\cref{ch:inference,ch:serving}.

\begin{figure}[t]
\centering
\resizebox{\linewidth}{!}{%
\begin{tikzpicture}[
  stage/.style={draw=sysblue,fill=sysblue!8,rounded corners=2pt,minimum width=2.25cm,
    minimum height=1.05cm,align=center,font=\small},
  sys/.style={draw=napkin,fill=napkin!12,rounded corners=2pt,minimum width=2.25cm,
    minimum height=0.8cm,align=center,font=\scriptsize},
  num/.style={circle,fill=sysblue,text=white,font=\scriptsize\bfseries,inner sep=1.5pt},
  arr/.style={-{Latex[length=2mm]},thick,sysblue!70}]
\node[stage] (s1) {Prepare\\dataset};
\node[stage,right=0.45cm of s1] (s2) {Build\\data loaders};
\node[stage,right=0.45cm of s2] (s3) {Load\\pretrained GPT};
\node[stage,right=0.45cm of s3] (s4) {Swap in a\\class head};
\node[stage,right=0.45cm of s4] (s5) {Loss and\\accuracy};
\node[stage,right=0.45cm of s5] (s6) {Fine-tune};
\node[stage,right=0.45cm of s6] (s7) {Classify new\\messages};
\foreach \i in {1,...,7}{\node[num] at (s\i.north west) {\i};}
\foreach \a/\b in {1/2,2/3,3/4,4/5,5/6,6/7}{\draw[arr] (s\a) -- (s\b);}
\node[sys,below=0.55cm of s2] (y2) {padding waste,\\length bucketing};
\node[sys,below=0.55cm of s4] (y4) {freezing: what is\\saved, what is not};
\node[sys,below=0.55cm of s6] (y6) {trunk feature\\caching};
\node[sys,below=0.55cm of s7] (y7) {prefill-only\\serving};
\foreach \i in {2,4,6,7}{\draw[densely dashed,napkin] (s\i) -- (y\i);}
\node[font=\scriptsize\itshape,text=napkin,anchor=east] at ([xshift=-0.2cm]y2.west) {systems layer};
\end{tikzpicture}}
\caption{Chapter roadmap. The top row is the modelling pipeline for turning a pretrained
GPT into a spam classifier; the dashed boxes are the systems questions we attach to each
stage.}
\label{fig:classify:roadmap}
\end{figure}

% =====================================================================================
\section{The fine-tuning landscape}\label{sec:classify:landscape}\covers{LLM 6.1}

Once pretraining (\cref{ch:pretrain}) is done, there are two classic ways to adapt a
language model to a task.

\begin{description}
\item[Instruction fine-tuning] trains the model to read a natural-language request and
answer in natural language. The output space stays the full vocabulary, so a single
model can handle a wide range of tasks: summarise, translate, fix the grammar, explain
a stack trace. The cost is that it needs diverse, carefully formatted data, and the
outputs are free text that must be parsed. This is \cref{ch:instruct}.
\item[Classification fine-tuning] trains the model to emit one label from a fixed set.
The output space shrinks to $C$ classes (spam/not-spam, sentiment, intent, toxicity
category). The model is useless for anything else, but it is cheap to train, its output
needs no parsing, and its accuracy is easy to measure.
\end{description}

In a production pipeline you often need both. An instruction-tuned assistant may sit
behind a battery of small classifiers that screen inputs, route requests to specialised
models, or detect policy violations. These classifiers run on \emph{every} request, so
their cost per token matters more than their absolute size suggests.

\begin{figure}[ht]
\centering
\begin{tikzpicture}[
  base/.style={draw=sysblue,fill=sysblue!15,rounded corners=3pt,minimum width=2.4cm,
    minimum height=0.9cm,align=center,font=\small\bfseries},
  leaf/.style={draw=#1,fill=#1!10,rounded corners=2pt,text width=4.2cm,align=center,
    font=\scriptsize,minimum height=0.95cm},
  arr/.style={-{Latex[length=2mm]},thick,srcgray}]
\node[base] (pt) at (0,0) {Pretrained\\base model};
\node[leaf=keygreen] (cls) at (4.6,1.35)
  {\textbf{Classification FT}\\$C$-way head, labelled data (this chapter)};
\node[leaf=sysblue] (ins) at (4.6,0)
  {\textbf{Instruction FT (SFT)}\\full vocabulary, prompt/response pairs (Ch.~\ref*{ch:instruct})};
\node[leaf=labpurple] (rl) at (4.6,-1.35)
  {\textbf{Preference / RL}\\reward on whole outputs (Ch.~\ref*{ch:grpo}--\ref*{ch:rlsys})};
\node[leaf=napkin,text width=3.0cm,minimum height=3.6cm] (peft) at (9.4,0)
  {\textbf{PEFT}\\(LoRA, QLoRA, adapters)\\[3pt]an orthogonal axis:\\\emph{which} parameters\\move
   (\S\ref*{sec:instruct:lora})};
\foreach \n in {cls,ins,rl}{\draw[arr] (pt.east) -- (\n.west);
  \draw[densely dashed,napkin,thick] (\n.east) -- (\n.east -| peft.west);}
\end{tikzpicture}
\caption{Ways to adapt a pretrained model. \emph{What} we optimise (labels, responses,
rewards) is one axis; \emph{which parameters} we update (all, the last few layers, or
low-rank adapters) is a separate axis that applies to every method.}
\label{fig:classify:landscape}
\end{figure}

A second axis, separate from the first, is \emph{how many parameters move}. You can update
all of them (full fine-tuning), only the top few layers plus the new head (the default
in this chapter), or small low-rank adapters (LoRA, \cref{sec:instruct:lora}). The
choice sets the memory footprint, the per-step compute, and how much capacity the model
has to adapt.

\begin{systemslens}[What each adaptation regime keeps in memory]
For a model with $N$ parameters trained with AdamW in mixed precision, the persistent
state is about $16N$ bytes (\cref{ch:pretrain}): bf16 weights (2), bf16 gradients (2),
fp32 master weights (4), and two fp32 Adam moments (8). Freezing a parameter removes
everything except the 2-byte weight. If only a fraction $f$ of the parameters is
trainable, the persistent state is roughly $2N + 14fN$ bytes. For GPT-2 small with only
the last block, the final LayerNorm and the head trainable, $f\approx 0.057$
(\cref{sec:classify:head}). The optimizer state then becomes a rounding error, and the
memory budget is set by weights and activations.
\end{systemslens}

\begin{napkinmath}[Full vs. partial fine-tuning of GPT-2 small]
GPT-2 small has $N\approx124\text{M}$ unique parameters. Full fine-tuning with AdamW:
$16 \times 124\text{M} \approx \SI{2.0}{GB}$ of persistent state. Training only the last
block, the final norm and a 2-way head ($\approx 7.1$M parameters):
$2 \times 124\text{M} + 14 \times 7.1\text{M} \approx 248 + 99 = \SI{347}{MB}$.
The six-fold reduction matters little on a datacenter GPU. On a laptop, or when one
GPU hosts dozens of task-specific classifiers, it decides whether the job fits.
\end{napkinmath}

% =====================================================================================
\section{Preparing the dataset}\label{sec:classify:data}\covers{LLM 6.2}

Our running example is the \emph{SMS Spam Collection}, a public corpus of 5{,}574 short
text messages. Each is labelled \code{ham} (legitimate, 4{,}827 messages) or
\code{spam} (747 messages). It is distributed as a tab-separated file, one
\code{label<TAB>text} pair per line. The data is small enough to fine-tune GPT-2 small
on a laptop in minutes, and realistic enough to show the usual problems: class
imbalance, long-tailed message lengths, and near-duplicate messages.

\paragraph{Class imbalance.} Only $747/5574 \approx 13.4\%$ of messages are spam. A model
that always predicts \code{ham} already scores 86.6\% accuracy, which makes accuracy a
poor signal early in training. The simplest fix is \emph{undersampling}: keep every
spam message and draw an equal-sized random subset of 747 ham messages, giving a
balanced dataset of 1{,}494 examples. Undersampling throws data away, which is
acceptable here because we want a fast, interpretable experiment. In production you
would more likely keep all the data and use one of the following:
\begin{itemize}
\item a class-weighted loss, where the weight of class $c$ is proportional to
$1/\text{freq}(c)$;
\item oversampling the minority class, either by repetition or by generating
paraphrases with an LLM;
\item reporting precision, recall and F1 (or the area under the precision--recall
curve) instead of accuracy, and choosing a decision threshold for the cost you care
about. Here that cost is usually a false positive: a real message hidden as spam.
\end{itemize}

\begin{lstlisting}[caption={Undersampling to balance classes, then a seeded three-way split (from \ident{llmsys/finetune.py}).},label={lst:classify:balance}]
def balance_by_undersampling(labels, seed=123):
    by_class = {}
    for i, y in enumerate(labels):
        by_class.setdefault(int(y), []).append(i)
    n_min = min(len(v) for v in by_class.values())
    rng = random.Random(seed)
    keep = []
    for y in sorted(by_class):
        idx = by_class[y]
        keep.extend(idx if len(idx) == n_min else rng.sample(idx, n_min))
    rng.shuffle(keep)
    return keep

def split_indices(n, train_frac=0.7, val_frac=0.1, seed=123):
    idx = list(range(n))
    random.Random(seed).shuffle(idx)
    a = int(n * train_frac)
    b = a + int(n * val_frac)
    return idx[:a], idx[a:b], idx[b:]       # test = the remainder
\end{lstlisting}

\paragraph{Splitting.} We shuffle once with a fixed seed and cut the balanced set
70/10/20 into training (1{,}045), validation (149) and test (300) examples. Labels are
mapped to integers, \code{ham}$\to 0$ and \code{spam}$\to 1$, because
\code{cross\_entropy} expects class indices.

\begin{figure}[ht]
\centering
\begin{tikzpicture}[x=0.0021cm,y=1cm]
% raw
\node[anchor=east,font=\small] at (-60,2.0) {raw};
\fill[sysblue!55] (0,1.75) rectangle (4827,2.25);
\fill[warnred!65] (4827,1.75) rectangle (5574,2.25);
\node[font=\scriptsize,white] at (2413,2.0) {ham 4{,}827};
\node[font=\scriptsize,anchor=west] at (5600,2.0) {spam 747};
% balanced
\node[anchor=east,font=\small] at (-60,1.0) {balanced};
\fill[sysblue!55] (0,0.75) rectangle (747,1.25);
\fill[warnred!65] (747,0.75) rectangle (1494,1.25);
\node[font=\scriptsize,anchor=west] at (1520,1.0) {747 + 747 = 1{,}494};
\draw[densely dashed,srcgray] (747,1.75) -- (747,1.25);
% split
\node[anchor=east,font=\small] at (-60,0.0) {split};
\fill[keygreen!55] (0,-0.25) rectangle (1045,0.25);
\fill[napkin!60] (1045,-0.25) rectangle (1194,0.25);
\fill[labpurple!55] (1194,-0.25) rectangle (1494,0.25);
\node[font=\scriptsize] at (522,0) {train 1{,}045};
\node[font=\scriptsize,anchor=south] at (1120,0.27) {val 149};
\node[font=\scriptsize,anchor=north] at (1344,-0.27) {test 300};
\draw[-{Latex[length=1.6mm]},srcgray] (2600,1.75) -- (1000,1.28);
\node[font=\scriptsize\itshape,text=srcgray,anchor=west] at (2650,1.45) {random subset of ham};
\end{tikzpicture}
\caption{From the raw SMS collection to a balanced, three-way split, drawn to scale.
Undersampling discards about 85\% of the ham messages.}
\label{fig:classify:split}
\end{figure}

\begin{pitfall}[Leakage through duplicates]
SMS corpora contain many exact and near-exact duplicates, especially among spam
campaigns that send the same template to thousands of numbers. If a message appears in
both training and test, test accuracy measures memorisation rather than generalisation.
Deduplicate \emph{before} splitting: hash a normalised form (lowercased, whitespace
collapsed, digits masked) and keep one copy per hash. The same discipline applies to the
much larger contamination problem in LLM evaluation (\cref{ch:eval}).
\end{pitfall}

% =====================================================================================
\section{Creating data loaders}\label{sec:classify:loaders}\covers{LLM 6.3}

Text messages vary in length, but a tensor is a rectangle. The SMS messages, tokenised
with the GPT-2 BPE tokenizer (\cref{ch:text}), run from a handful of tokens to well over
a hundred. There are two simple ways to make them rectangular:

\begin{enumerate}
\item \textbf{Truncate} every message to the length of the shortest one. This is cheap
but destroys information, because most spam cues sit beyond the first few tokens.
\item \textbf{Pad} every message with a filler token up to some length $T_{\max}$. Since
GPT-2 has no dedicated padding token, we reuse \code{<|endoftext|>} (id 50256). That is
safe as long as the model never reads the hidden states at padded positions; see
\cref{sec:classify:head}.
\end{enumerate}

The simplest padding policy, and the one used in the LLM book, fixes $T_{\max}$ to the
longest training message (120 tokens for this dataset) and pads the validation and test
sets to the same length, truncating anything longer. Every batch then has the same
shape, which is convenient for debugging and for compilers that specialise on shapes.

\begin{figure}[ht]
\centering
\begin{tikzpicture}[
  tok/.style={draw=sysblue!60,fill=sysblue!12,minimum width=0.95cm,minimum height=0.5cm,
    font=\scriptsize\ttfamily,inner sep=1pt},
  pad/.style={tok,draw=srcgray!60,fill=srcgray!12,text=srcgray},
  lab/.style={font=\scriptsize,anchor=east,text width=3.5cm,align=right}]
\node[lab] at (-0.2,0) {``Call now to claim''};
\node[lab] at (-0.2,-0.65) {``ok see you''};
\node[lab] at (-0.2,-1.3) {``Free entry in a weekly''};
\foreach \x/\t in {0/14879,1/783,2/284,3/1624}{\node[tok] at (\x*1.0+0.5,0) {\t};}
\foreach \x in {4,5}{\node[pad] at (\x*1.0+0.5,0) {50256};}
\foreach \x/\t in {0/482,1/766,2/345}{\node[tok] at (\x*1.0+0.5,-0.65) {\t};}
\foreach \x in {3,4,5}{\node[pad] at (\x*1.0+0.5,-0.65) {50256};}
\foreach \x/\t in {0/11146,1/5726,2/287,3/257,4/10273,5/1366}{\node[tok] at (\x*1.0+0.5,-1.3) {\t};}
\draw[decorate,decoration={brace,amplitude=4pt,mirror}] (0,-1.7) -- (6,-1.7)
  node[midway,below=4pt,font=\scriptsize] {$T_{\max}$ (longest sequence)};
\node[font=\scriptsize,text=srcgray,anchor=west,text width=3.2cm] at (6.3,-0.4)
  {grey cells: \code{<|endoftext|>} used as padding; real ids shown are illustrative};
\end{tikzpicture}
\caption{Tokenise, then pad to a common length. Every grey cell costs as much compute
as a real token and, in attention, adds a row and a column to the score matrix.}
\label{fig:classify:pad}
\end{figure}

\paragraph{Dynamic padding.} A better default is to pad each \emph{batch} to its own
longest member. The dataset stores variable-length token lists, and the collate function
pads them when the batch is assembled. It also returns the true lengths, which we will
need to find the last real token.

\begin{lstlisting}[caption={A dataset that stores token lists, and a collate function that pads per batch and returns lengths.},label={lst:classify:collate}]
class ClassificationDataset(Dataset):
    def __init__(self, texts, labels, tokenizer, max_length=None):
        self.encoded = [tokenizer.encode(t)[:max_length] if max_length
                        else tokenizer.encode(t) for t in texts]
        self.labels = [int(y) for y in labels]
    def __len__(self):
        return len(self.labels)
    def __getitem__(self, i):
        return self.encoded[i], self.labels[i]

def pad_batch(seqs, pad_id, side="right", max_length=None):
    if max_length is not None:
        seqs = [list(s)[:max_length] for s in seqs]
    T = max(len(s) for s in seqs)
    ids = torch.full((len(seqs), T), pad_id, dtype=torch.long)
    for i, s in enumerate(seqs):
        t = torch.as_tensor(list(s), dtype=torch.long)
        if side == "right": ids[i, :len(s)] = t
        else:               ids[i, T - len(s):] = t
    return ids, torch.tensor([len(s) for s in seqs])

def classification_collate(batch, pad_id, side="right"):
    seqs, labels = zip(*batch)
    ids, lengths = pad_batch(seqs, pad_id, side)
    return ids, lengths, torch.tensor(labels)
\end{lstlisting}

We wrap the three splits in \code{DataLoader}s with batch size 8, shuffling and
\code{drop\_last=True} for training only. Dropping the last incomplete batch keeps every
optimisation step at the same batch size. Validation and test keep all their examples,
so every example is scored. A quick sanity check is to iterate once and print the
shapes: with batch size 8 and global padding, input batches are $8 \times 120$ and label
batches are $8$. With dynamic padding the second dimension varies from batch to batch.

\begin{systemslens}[The real price of padding]
Padding costs more than the extra tokens suggest. For a batch of $B$ sequences padded to
$T$, the MLP and projection FLOPs scale with $BT$, while attention scores scale with
$BT^2$. A single 120-token outlier in a batch of 10-token messages makes that batch
about $12\times$ more expensive in the linear layers and about $144\times$ more
expensive in the attention score matrix. Two levers fix this:
\begin{itemize}
\item \textbf{Length bucketing.} Group sequences of similar length into the same batch.
Sorting the whole dataset by length destroys the randomness SGD relies on. A standard
compromise is to shuffle, cut the data into \emph{mega-batches} of, say,
$50\times$ the batch size, sort within each mega-batch, slice into batches, and shuffle
the batch order (\cref{lst:classify:bucket}).
\item \textbf{Token-budget batching.} Fix the number of \emph{tokens} per batch
instead of the number of sequences, so batches of short messages hold more of them.
This keeps GPU utilisation and memory use flat. Large-scale training and serving
systems almost always batch by tokens.
\end{itemize}
The next step after bucketing is \emph{packing} (\cref{sec:instruct:packing}), which
removes padding completely by concatenating sequences and masking attention between
them.
\end{systemslens}

\begin{lstlisting}[caption={Sort-within-mega-batch bucketing and a padding-waste meter.},label={lst:classify:bucket}]
def length_bucketed_batches(lengths, batch_size, seed=0, bucket_mult=50):
    rng = random.Random(seed)
    idx = list(range(len(lengths)))
    rng.shuffle(idx)
    mega = batch_size * bucket_mult
    batches = []
    for s in range(0, len(idx), mega):
        chunk = sorted(idx[s:s + mega], key=lambda i: lengths[i])
        batches.extend(chunk[j:j + batch_size] for j in range(0, len(chunk), batch_size))
    rng.shuffle(batches)                  # randomise batch order
    return batches

def padding_stats(lengths, batches):
    real = slots = 0
    for b in batches:
        L = [lengths[i] for i in b]
        real += sum(L)
        slots += max(L) * len(L)
    return dict(real_tokens=real, padded_slots=slots, waste=1 - real / max(slots, 1))
\end{lstlisting}

\Cref{fig:classify:waste} shows what bucketing buys. We generated 4{,}000 synthetic
message lengths with a log-normal distribution (median 13 tokens, maximum 108), similar
in shape to real SMS data, and measured the fraction of computed token slots that are
padding. Padding everything to the global maximum wastes 86\% of all slots. Dynamic
padding with random batches wastes 40--76\%, rising with batch size because a larger
batch is more likely to contain an outlier. Bucketing keeps waste below 10\% at every
batch size.

\begin{figure}[ht]
\centering
\begin{tikzpicture}
\begin{axis}[width=0.78\linewidth,height=5.2cm,xmode=log,log basis x=2,
  xlabel={batch size},ylabel={padding waste (fraction of slots)},
  ymin=0,ymax=0.9,xtick={4,8,16,32,64,128},xticklabels={4,8,16,32,64,128},
  legend style={font=\scriptsize,at={(0.98,0.42)},anchor=east},grid=major,
  grid style={gray!20}]
\addplot[mark=*,thick,warnred] coordinates {(4,0.416)(8,0.528)(16,0.61)(32,0.667)(64,0.721)(128,0.761)};
\addlegendentry{random batches, dynamic padding}
\addplot[mark=square*,thick,keygreen] coordinates {(4,0.04)(8,0.049)(16,0.062)(32,0.088)(64,0.099)(128,0.077)};
\addlegendentry{length-bucketed batches}
\addplot[dashed,thick,srcgray,domain=4:128] {0.856};
\addlegendentry{pad all to global max}
\end{axis}
\end{tikzpicture}
\caption{Padding waste versus batch size, measured in this environment on 4{,}000
synthetic log-normal message lengths (\ident{labs/lab16_classify}). Bucketing with
mega-batches of $50\times$ the batch size removes nearly all of the waste.}
\label{fig:classify:waste}
\end{figure}

% =====================================================================================
\section{Initializing a model with pretrained weights}\label{sec:classify:init}\covers{LLM 6.4}

We start from GPT-2 small (124M) and load the released weights into the
\code{GPTModel} of \cref{ch:gpt} with \code{load\_gpt2\_hf\_state\_dict}
(\cref{ch:pretrain}). Two settings must match the checkpoint: the released model has
biases on the query/key/value projections (\code{qkv\_bias=True}), and it uses a
1{,}024-token context. We also set dropout to zero, since we will not fine-tune long
enough for regularisation to matter.

\begin{lstlisting}[caption={Loading GPT-2 small for fine-tuning.},label={lst:classify:load}]
from safetensors.torch import load_file
from llmsys.gpt import GPTModel, gpt2_config, load_gpt2_hf_state_dict

cfg = gpt2_config("gpt2-small (124M)", drop_rate=0.0)   # qkv_bias=True, ctx 1024
model = GPTModel(cfg)
load_gpt2_hf_state_dict(model, load_file("gpt2/model.safetensors"))
model.eval()
\end{lstlisting}

Before changing anything, run two sanity checks. First, generate a continuation of
a neutral prompt (``Every effort moves you'') with greedy decoding and confirm the text
is fluent English. A garbled continuation means the weight mapping is wrong, most often
because a Conv1D weight was not transposed. Second, see whether the base model can
already classify when prompted:

\begin{lstlisting}[style=plain]
Is the following text 'spam'? Answer with 'yes' or 'no': 'You are a winner you
have been specially selected to receive $1000 cash or a $2000 award.'
\end{lstlisting}

A base GPT-2 typically repeats the question or wanders off instead of answering. It was
never trained to follow instructions (\cref{ch:instruct}). That is the motivation for
replacing the output layer and training it directly.

\begin{keyidea}
Fine-tuning works because the pretrained trunk already produces hidden states that
encode much of what a classifier needs: topic, register, sentiment, and the
presence of URLs, prize language or urgency cues. The new head only has to find a
linear boundary in that representation space. Unfreezing the last block lets the
features themselves bend slightly toward the task.
\end{keyidea}


\subsection{Where does the task information live?}\label{sec:classify:probe}

Before deciding how many blocks to unfreeze, it helps to measure how linearly
separable the classes already are at each depth of the frozen network. A \emph{linear
probe} trains only a logistic-regression layer on the last-real-token hidden state taken
after block $\ell$, for each $\ell$. Forward hooks collect all twelve representations
in a single pass over the data:

\begin{lstlisting}[caption={Collecting per-layer last-token features with forward hooks.},label={lst:classify:probe}]
@torch.no_grad()
def collect_layer_features(model, loader, device="cpu"):
    feats = {i: [] for i in range(len(model.trf_blocks))}
    current = {}
    hooks = [blk.register_forward_hook(
                 lambda m, inp, out, i=i: current.__setitem__(i, out))
             for i, blk in enumerate(model.trf_blocks)]
    labels = []
    for ids, lengths, y in loader:
        model(ids.to(device))
        rows = torch.arange(ids.shape[0])
        for i, h in current.items():
            feats[i].append(h[rows, lengths - 1].float().cpu())  # (B, d)
        labels.append(y)
    for h in hooks:
        h.remove()
    return {i: torch.cat(v) for i, v in feats.items()}, torch.cat(labels)
\end{lstlisting}

The resulting curve of probe accuracy against depth is a cheap guide to where to cut the
network. If accuracy is already high at block 8, you can drop blocks 9--12 entirely for
this task, which cuts serving cost by a third. Or you can freeze through block 8 and
cache its outputs (\cref{sec:classify:head}). The memory cost of probing is the number
of examples $\times$ layers $\times$ $d$ floats: $1{,}045\times12\times768\times4$
bytes $\approx\SI{39}{MB}$, small because only one position per message is kept.
\Cref{ex:classify:probe} asks you to produce the curve for GPT-2 small on the SMS data.

% =====================================================================================
\section{Adding a classification head}\label{sec:classify:head}\covers{LLM 6.5}

The pretrained \code{out\_head} maps the final 768-dimensional hidden state to 50{,}257
vocabulary logits. We replace it with \code{nn.Linear(768, 2)}: one logit for \code{ham}
and one for \code{spam}. With two classes we could get away with a single logit and a
sigmoid. Keeping $C$ logits and the softmax cross-entropy loss means the same code
handles any number of classes.

\paragraph{Which layers to train.} Lower transformer layers in pretrained LMs tend to
encode general lexical and syntactic features that transfer across tasks. Upper layers
are more specialised to next-token prediction. A common, cheap recipe is therefore:

\begin{enumerate}
\item freeze every parameter;
\item replace \code{out\_head} (new parameters are trainable by default);
\item unfreeze the last transformer block and the final LayerNorm.
\end{enumerate}

Training more blocks usually helps a little and costs more; training all of them is
full fine-tuning. \Cref{ex:classify:layers} asks you to measure the trade-off.

\begin{lstlisting}[caption={Swapping the head and choosing what to train.},label={lst:classify:head}]
def add_classification_head(model, num_classes, trainable_blocks=1,
                            train_final_norm=True):
    for p in model.parameters():
        p.requires_grad = False
    emb_dim = model.out_head.in_features
    model.out_head = nn.Linear(emb_dim, num_classes)     # trainable
    if trainable_blocks > 0:
        for blk in model.trf_blocks[-trainable_blocks:]:
            for p in blk.parameters():
                p.requires_grad = True
    if train_final_norm:
        for p in model.final_norm.parameters():
            p.requires_grad = True
    return model
\end{lstlisting}

\begin{figure}[ht]
\centering
\begin{tikzpicture}[
  blk/.style={draw=#1,fill=#1!12,minimum width=5.2cm,minimum height=0.55cm,font=\scriptsize,
    rounded corners=1.5pt},
  arr/.style={-{Latex[length=1.6mm]},srcgray}]
\node[blk=srcgray] (emb) at (0,0) {token + position embeddings \quad(frozen)};
\node[blk=srcgray] (b1) at (0,0.8) {transformer block 1 \quad(frozen)};
\node[font=\small] (dots) at (0,1.45) {$\vdots$};
\node[blk=srcgray] (b11) at (0,2.1) {transformer block 11 \quad(frozen)};
\node[blk=keygreen] (b12) at (0,2.9) {transformer block 12 \quad(trainable)};
\node[blk=keygreen] (ln) at (0,3.7) {final LayerNorm \quad(trainable)};
\node[blk=warnred,minimum width=2.4cm] (head) at (0,4.6) {Linear $768\to 2$ (new)};
\node[draw=srcgray,dashed,minimum width=3.4cm,minimum height=0.55cm,font=\scriptsize,
  text=srcgray] (old) at (5.0,4.6) {Linear $768\to 50257$};
\draw[arr] (emb)--(b1); \draw[arr] (b11)--(b12); \draw[arr] (b12)--(ln); \draw[arr] (ln)--(head);
\draw[thick,warnred] ($(old.north west)$) -- ($(old.south east)$);
\draw[-{Latex[length=1.6mm]},warnred,thick] (old.west) -- (head.east)
  node[midway,above,font=\scriptsize] {replace};
\node[font=\scriptsize,align=left,anchor=west] at (3.0,2.5)
  {\textcolor{keygreen}{\rule{6pt}{6pt}} gradients + AdamW state\\
   \textcolor{srcgray}{\rule{6pt}{6pt}} weights only, no graph kept};
\end{tikzpicture}
\caption{The classifier: a frozen GPT-2 trunk, a trainable top block and final norm,
and a new two-way head in place of the vocabulary projection.}
\label{fig:classify:frozen}
\end{figure}

With this recipe, 7{,}090{,}946 of the model's parameters are trainable: the last block
(7{,}087{,}872), the final norm (1{,}536) and the head ($768\times2+2 = 1{,}538$). That is
5.7\% of the network.

\begin{systemslens}[Freezing saves backward compute and activation memory, not just optimizer state]
Autograd records a node only for operations whose inputs require gradients. With blocks
1--11 frozen, their inputs (the embedding outputs) do not require gradients, so no
graph is recorded for them and their intermediate activations are freed right after
the forward pass. Backward stops at the input of block 12. Two consequences:
\begin{itemize}
\item \textbf{Compute.} A training step costs roughly $2N$ FLOPs per token forward plus
$4N_{\text{trainable}}$ backward, instead of $6N$. With one of twelve blocks trainable,
that is about $(2 + 4/12)/6 \approx 0.39$ of a full fine-tuning step, or $2.6\times$
cheaper.
\item \textbf{Trunk caching.} If the bottom 11 blocks never change, their output for a
given message never changes either. You can run them once over the dataset, store the
$T\times 768$ hidden states, and train only block 12 and the head on the cache. The
stored hidden states fit easily: 1{,}045 messages $\times$ 120 positions
$\times$ 768 dimensions $\times$ 2 bytes $\approx \SI{193}{MB}$. This is a big win when
you sweep hyperparameters over many epochs.
\end{itemize}
LoRA (\cref{sec:instruct:lora}) is a different case. Adapters sit in \emph{every}
layer, so the activation gradient must flow through the whole network, and the savings
come only from skipping the weight-gradient GEMMs of the frozen weights.
\end{systemslens}

\subsection{Which token should the head read?}

The model returns logits of shape $(B, T, C)$, one $C$-vector per position. We need one
vector per message. Recall the causal mask from \cref{ch:attention}: position $t$ can
attend only to positions $0,\dots,t$. Only the \emph{last} position has attended to the
whole message, so its hidden state is the natural summary to classify. The first
position has seen a single token.

\begin{figure}[ht]
\centering
\begin{tikzpicture}[scale=0.55]
\foreach \i in {0,...,5}{
  \foreach \j in {0,...,5}{
    \pgfmathsetmacro{\allowed}{\j<=\i ? 1 : 0}
    \ifnum\j>\i
      \fill[srcgray!15] (\j,-\i) rectangle ++(1,-1);
    \else
      \ifnum\i=5
        \fill[keygreen!55] (\j,-\i) rectangle ++(1,-1);
      \else
        \fill[sysblue!25] (\j,-\i) rectangle ++(1,-1);
      \fi
    \fi
    \draw[white,thick] (\j,-\i) rectangle ++(1,-1);
  }}
\foreach \w [count=\k from 0] in {Free,entry,in,a,weekly,draw}{
  \node[font=\scriptsize,anchor=east] at (-0.1,-\k-0.5) {\w};
  \node[font=\scriptsize,rotate=55,anchor=west] at (\k+0.35,0.15) {\w};}
\draw[thick,keygreen!70!black,-{Latex[length=2mm]}] (6.3,-5.5) -- (8.5,-5.5)
  node[right,font=\scriptsize,align=left] {last row sees every token\\$\Rightarrow$ feed this
  hidden state to the head};
\node[font=\scriptsize,anchor=west,text=srcgray] at (6.3,-1) {grey: masked (future)};
\end{tikzpicture}
\caption{Causal attention over a six-token message. Only the final position aggregates
information from the whole sequence, so its output is the one we classify.}
\label{fig:classify:lasttoken}
\end{figure}

With \emph{right} padding, the last position of a padded row is a padding token, not the
last real token. The LLM book reads \code{logits[:, -1, :]} anyway. That works
reasonably well, because the padding token at the end still attends to every real
token before it. But it makes the prediction depend on \emph{how much} padding follows
the message, which in turn depends on which other messages happen to share the batch.
Two cleaner options:

\begin{enumerate}
\item \textbf{Gather the last real token.} With the lengths returned by the collate
function, pick row $b$ at position $\ell_b - 1$. The prediction is then exactly
invariant to padding (the lab tests this).
\item \textbf{Left padding with a key-padding mask.} Put the padding before the message
and mask padded keys (the \code{pad\_mask} argument of the core models). The last
position is now always real, so \code{logits[:, -1]} is correct. This is what batched
generation does (\cref{ch:serving}). With learned absolute positions (GPT-2) the real
tokens shift to higher position ids, which changes the outputs slightly; with RoPE it
does not.
\end{enumerate}

\begin{lstlisting}[caption={Selecting the per-message logits.},label={lst:classify:select}]
def select_class_logits(logits, lengths=None, side="right"):
    """logits: (B, T, C) -> (B, C)."""
    if lengths is None or side == "left":
        return logits[:, -1, :]
    idx = (lengths - 1).clamp(min=0).to(logits.device)
    return logits[torch.arange(logits.shape[0], device=logits.device), idx]
\end{lstlisting}

\begin{deepdive}[Other pooling choices]
Encoder models such as BERT classify from a dedicated \code{[CLS]} token placed at the
\emph{start}. That works only because encoders use bidirectional attention. In a causal
decoder you can also mean-pool the hidden states of the real tokens. This is sometimes
more robust for long inputs, but it dilutes the strongest signal, which tends to sit at
the end. Embedding models built from decoders (e5-mistral-style) typically use
last-token pooling with an appended end-of-sequence token, which is the recipe of this
chapter.
\end{deepdive}

% =====================================================================================
\section{Calculating the classification loss and accuracy}\label{sec:classify:loss}\covers{LLM 6.6}

For a message $x$ with selected logits $\vect{z}\in\R^{C}$, the predicted class
distribution is $p(c\mid x) = \softmax(\vect{z})_c$. The predicted label is
$\hat{y} = \argmax_c z_c$. Softmax is monotonic, so we can take the $\argmax$ of the
logits directly. Training minimises the mean cross-entropy over a batch,
\begin{equation}
\mathcal{L} = -\frac{1}{B}\sum_{b=1}^{B} \log \softmax(\vect{z}_b)_{y_b}
            = \frac{1}{B}\sum_{b=1}^{B}\Bigl(\log\sum_{c} e^{z_{b,c}} - z_{b,y_b}\Bigr),
\label{eq:classify:ce}
\end{equation}
the same loss as in pretraining (\cref{ch:pretrain}). It is applied once per message
instead of once per token, over $C=2$ classes instead of 50{,}257. We cannot optimise
accuracy directly because it is piecewise constant, so its gradient is zero almost
everywhere. Cross-entropy is a smooth surrogate for it.

\begin{lstlisting}[caption={Loss on a batch and accuracy over a loader.},label={lst:classify:acc}]
def classification_loss(model, ids, lengths, labels, device="cpu", use_lengths=True):
    ids, labels = ids.to(device), labels.to(device)
    logits = select_class_logits(model(ids), lengths if use_lengths else None)
    return F.cross_entropy(logits.float(), labels)

@torch.no_grad()
def accuracy_loader(model, loader, device="cpu", num_batches=None, use_lengths=True):
    model.eval()
    correct = total = 0
    for i, (ids, lengths, labels) in enumerate(loader):
        if num_batches is not None and i >= num_batches:
            break
        logits = select_class_logits(model(ids.to(device)),
                                     lengths if use_lengths else None)
        correct += (logits.argmax(-1).cpu() == labels).sum().item()
        total += labels.numel()
    return correct / max(total, 1)
\end{lstlisting}

Evaluated before any fine-tuning, the new head is random, so accuracy on the balanced
splits sits near 50\%, which is chance. The \code{num\_batches} argument estimates
accuracy from a few batches, which is useful for frequent cheap checks during training.

\begin{pitfall}[Computing the loss in low precision]
Under bf16 autocast the logits come out in bf16, which has only 8 bits of mantissa. A
log-sum-exp over large vocabularies loses accuracy in that format. With two classes
the damage is small, but the habit of casting logits to fp32 (\code{logits.float()})
before \code{cross\_entropy} costs nothing here and matters a lot in
\cref{ch:pretrain,ch:grpo}.
\end{pitfall}

% =====================================================================================
\section{Fine-tuning the model on supervised data}\label{sec:classify:train}\covers{LLM 6.7}

The training loop is the pretraining loop of \cref{ch:pretrain} with three changes: the
loss is computed on the selected logits, only trainable parameters are passed to the
optimizer, and we track accuracy alongside loss. Following the LLM book we use AdamW
with learning rate $5\times10^{-5}$ and weight decay $0.1$ for about five epochs. On a
modern laptop CPU this takes several minutes for GPT-2 small; on a GPU, well under a
minute.

\begin{lstlisting}[caption={Classification fine-tuning loop (lab reference).},label={lst:classify:train}]
def train_classifier(model, train_loader, val_loader, epochs=3, lr=5e-4,
                     weight_decay=0.1, device="cpu", use_lengths=True, log=print):
    params = [p for p in model.parameters() if p.requires_grad]
    opt = torch.optim.AdamW(params, lr=lr, weight_decay=weight_decay)
    history = []
    for ep in range(epochs):
        model.train()
        total, n, correct, seen = 0.0, 0, 0, 0
        for ids, lengths, labels in train_loader:
            opt.zero_grad(set_to_none=True)
            logits = select_class_logits(model(ids.to(device)),
                                         lengths if use_lengths else None)
            loss = F.cross_entropy(logits.float(), labels.to(device))
            loss.backward()
            opt.step()
            total += loss.item(); n += 1
            correct += (logits.argmax(-1).cpu() == labels).sum().item()
            seen += len(labels)
        va_acc = accuracy_loader(model, val_loader, device, use_lengths=use_lengths)
        history.append(dict(epoch=ep + 1, loss=total / n,
                            train_acc=correct / seen, val_acc=va_acc))
    return history
\end{lstlisting}

Two details are worth noticing. First, the training accuracy is accumulated from
forward passes we already pay for, instead of a second pass over the training set. That
pass would cost as much as a full extra epoch of forward compute. The number is
slightly stale (it mixes predictions from the start and end of the epoch), which is
fine for monitoring. Second, \code{loss.item()} forces a host--device synchronisation
on every step. On a GPU, it is better to accumulate the loss in a device tensor and
read it once per logging interval (\cref{ch:pytorch}).

\Cref{fig:classify:curves} shows a training run on the lab's synthetic spam task: a
2-layer, 64-dimensional GPT trained from scratch, measured on CPU in this environment.
The shape is what you will see with real GPT-2 on real SMS data. The loss falls
quickly, training accuracy saturates near 100\%, and validation accuracy plateaus a few
points lower. With a pretrained GPT-2 on the SMS data, the LLM book reports training
and validation accuracies in the high 90s and test accuracy in the mid 90s.

\begin{figure}[ht]
\centering
\begin{tikzpicture}
\begin{groupplot}[group style={group size=2 by 1,horizontal sep=1.6cm},
  width=0.46\linewidth,height=4.6cm,xlabel={epoch},xtick={1,2,3,4},grid=major,
  grid style={gray!20},legend style={font=\scriptsize}]
\nextgroupplot[ylabel={training loss},ymin=0]
\addplot[mark=*,thick,sysblue] coordinates {(1,0.680)(2,0.252)(3,0.138)(4,0.087)};
\nextgroupplot[ylabel={accuracy},ymin=0.5,ymax=1.02,legend pos=south east]
\addplot[mark=*,thick,sysblue] coordinates {(1,0.805)(2,0.973)(3,0.984)(4,0.995)};
\addlegendentry{train}
\addplot[mark=square*,thick,napkin] coordinates {(1,0.635)(2,0.904)(3,0.904)(4,0.942)};
\addlegendentry{validation}
\end{groupplot}
\end{tikzpicture}
\caption{Loss and accuracy per epoch on the synthetic spam task of Lab~16 (600 messages,
40\% spam before balancing, tiny GPT, AdamW at $2\times10^{-3}$, length-bucketed batches
of 16). Measured on CPU in this environment.}
\label{fig:classify:curves}
\end{figure}

\paragraph{Reading the curves.} If validation accuracy lags training accuracy by more
than a few points and the gap keeps widening, the model is memorising. The usual fixes,
in rough order of cost, are fewer epochs (early stopping on validation accuracy),
higher weight decay, dropout, more data, and training fewer layers. If both curves
stall near 50\%, the usual causes are (i) the head is reading padding positions while
the padding amount is correlated with the label, (ii) the learning rate is too low for
a freshly initialised head, or (iii) the trainable flags are wrong. Print
\code{sum(p.numel() for p in model.parameters() if p.requires\_grad)} to check.

\begin{napkinmath}[How long should this take on a GPU?]
Each training message costs about $2N + 4N_{\text{tr}}$ FLOPs per token
(\cref{sec:classify:head}), where $N$ counts only the non-embedding weights that do
work. For GPT-2 small, $N\approx 85$M and $N_{\text{tr}}\approx 7.1$M, giving
$\approx 1.98\times10^8$ FLOPs per token. Five epochs over 1{,}045 messages padded to
120 tokens is $6.3\times10^5$ token-positions, or about $1.2\times10^{14}$ FLOPs. At a
pessimistic 20\% of an A100's \SI{312}{TFLOP/s}, that is two seconds of arithmetic. In
practice a run takes much longer: kernel launches for small batches, Python overhead
and evaluation passes dominate. This chapter's workload is \emph{latency-bound}, not
FLOP-bound. Larger batches, \code{torch.compile} and CUDA graphs (\cref{ch:inference})
are the relevant levers.
\end{napkinmath}

% =====================================================================================
\section{Using the LLM as a spam classifier}\label{sec:classify:use}\covers{LLM 6.8}

To classify a new message we tokenise it, truncate it to the model's context
length (or to the training $T_{\max}$, if we trained with global padding and want
identical conditions), run one forward pass, and take the $\argmax$ of the selected
logits.

\begin{lstlisting}[caption={Inference with the fine-tuned classifier.},label={lst:classify:infer}]
@torch.no_grad()
def classify(model, tokenizer, text, device="cpu", max_length=None,
             labels=("not spam", "spam")):
    ctx = model.cfg["context_length"]
    ids = tokenizer.encode(text)[: min(max_length or ctx, ctx)]
    logits = model(torch.tensor([ids], device=device))[:, -1, :]   # no padding
    return labels[int(logits.argmax(-1))]

classify(model, tok, "You are a winner! Reply now to claim your free prize.")
# -> 'spam'
classify(model, tok, "Are we still on for dinner tonight?")
# -> 'not spam'

torch.save(model.state_dict(), "spam_classifier.pth")   # 2-way head included
\end{lstlisting}

To reload, rebuild the architecture \emph{with} the 2-way head before calling
\code{load\_state\_dict}. Otherwise the stored \code{out\_head.weight} of shape
$(2, 768)$ will not fit a $(50257, 768)$ parameter. Storing the full model costs about
$\SI{0.5}{GB}$ in fp32. If many tasks share one frozen trunk, store only the
trainable parameters per task ($\approx$\SI{28}{MB} here) together with a pointer to
the base checkpoint, which is the same idea as LoRA adapters in \cref{ch:instruct}.

\begin{systemslens}[Serving a classifier is a pure-prefill workload]
Generative serving alternates a compute-bound \emph{prefill} with a memory-bound
\emph{decode} loop (\cref{ch:inference}). A classifier has only the prefill. Every
request is a single forward pass over its own tokens, the output is $C$ numbers, and
the KV cache is discarded immediately. Consequences:
\begin{itemize}
\item Throughput is governed by FLOPs, not memory bandwidth, so large batches and
tensor-core-friendly shapes pay off directly.
\item The vocabulary head disappears. In GPT-2 small the $768\times50257$ projection
is 38.6M parameters, nearly half the size of all twelve blocks. Removing it saves those
FLOPs for every request.
\item Batching policy is all about padding. Bucket by length, or use a varlen
attention kernel that accepts a packed batch with per-sequence offsets
(\cref{sec:instruct:packing}).
\item Latency is easy to bound: there is no generation length to predict, so SLOs can
be met with simple admission control.
\end{itemize}
\end{systemslens}

\begin{napkinmath}[Classifier throughput on one H100]
Forward FLOPs per token for GPT-2 small without the vocabulary head:
$2\times 85\text{M} = 1.7\times10^8$ (attention terms are negligible at 30 tokens). At
50\% of the H100's \SI{989}{TFLOP/s}:
$4.9\times10^{14} / 1.7\times10^8 \approx 2.9$ million tokens per second. At about 30
tokens per SMS, that is roughly 97{,}000 messages per second from one GPU, if batches are
large and padding-free. With random batching at 60\% padding waste
(\cref{fig:classify:waste}), the same GPU delivers about 40\% of that.
\end{napkinmath}

\begin{deepdive}[When not to fine-tune a 124M-parameter decoder]
For a binary task with a few thousand examples, a linear probe on frozen sentence
embeddings, or a small encoder (DistilBERT/ModernBERT-class), is often as accurate and
10--100$\times$ cheaper to serve. A decoder LLM classifier makes most sense when (i) you
already serve that trunk and can share it, (ii) the inputs are long or need world
knowledge, or (iii) you will later turn the same model into an instruction follower. An
LLM can also act as a \emph{zero-shot} classifier by scoring the log-probabilities of
label words. \Cref{ch:refine} develops this kind of probability-based scoring.
\end{deepdive}

% =====================================================================================
\section{From one classifier to a fleet}\label{sec:classify:fleet}

The LLM book stops once the spam classifier works. In production a single classifier is
rarely the whole story: you need to decide where to put the decision threshold, you may
run several classifiers on every request, and you will want an expensive model to see
only the inputs the cheap one is unsure about. This section collects the systems
patterns that follow from that.

\subsection{Metrics, thresholds and calibration}

The $\argmax$ rule implicitly thresholds $p(\text{spam}\mid x)$ at $0.5$. The costs of
the two error types are rarely equal. Hiding a real message from a user (a false
positive) is usually worse than letting a spam through (a false negative). With
$\mathrm{TP},\mathrm{FP},\mathrm{FN},\mathrm{TN}$ counted at a threshold $\tau$,
\begin{equation}
\text{precision}(\tau) = \frac{\mathrm{TP}}{\mathrm{TP}+\mathrm{FP}},\qquad
\text{recall}(\tau) = \frac{\mathrm{TP}}{\mathrm{TP}+\mathrm{FN}},\qquad
F_1 = \frac{2\,\text{precision}\cdot\text{recall}}{\text{precision}+\text{recall}}.
\label{eq:classify:prf}
\end{equation}
Raising $\tau$ trades recall for precision. The operating point is chosen on the
\emph{validation} split, for example the smallest $\tau$ whose validation precision is
at least 99\%, and then reported once on the test split.

\begin{figure}[ht]
\centering
\begin{tikzpicture}[cell/.style={minimum width=2.9cm,text width=2.6cm,minimum height=1.1cm,
    align=center,font=\scriptsize,draw=white,line width=1.5pt}]
\node[cell,fill=keygreen!35] at (0,0) {\textbf{TP}\\spam caught};
\node[cell,fill=warnred!35] at (2.9,0) {\textbf{FN}\\spam delivered (annoying)};
\node[cell,fill=warnred!65] at (0,-1.1) {\textbf{FP}\\real message hidden (costly)};
\node[cell,fill=keygreen!20] at (2.9,-1.1) {\textbf{TN}\\real message delivered};
\node[font=\scriptsize\bfseries] at (1.45,1.15) {predicted};
\node[font=\scriptsize] at (0,0.78) {spam}; \node[font=\scriptsize] at (2.9,0.78) {ham};
\node[font=\scriptsize\bfseries,rotate=90] at (-2.05,-0.55) {actual};
\node[font=\scriptsize,rotate=90] at (-1.7,0) {spam};
\node[font=\scriptsize,rotate=90] at (-1.7,-1.1) {ham};
\draw[-{Latex[length=2mm]},thick,sysblue] (4.7,0.3) -- (4.7,-1.4)
  node[midway,right,font=\scriptsize,align=left] {raise $\tau$:\\fewer FP,\\more FN};
\end{tikzpicture}
\caption{The confusion matrix with the predicted class as columns. The two error cells
cost different amounts, so the threshold $\tau$ on $p(\text{spam}\mid x)$ is a product
decision, not a modelling default.}
\label{fig:classify:confusion}
\end{figure}

A threshold only means something if the probabilities are \emph{calibrated}: among
messages scored $0.9$, about 90\% should really be spam. Fine-tuned networks are often
over-confident. A standard check is the expected calibration error. Bin the $n$ test
predictions by confidence into bins $B_1,\dots,B_M$ and compute
\begin{equation}
\mathrm{ECE} = \sum_{m=1}^{M}\frac{|B_m|}{n}\,\bigl|\mathrm{acc}(B_m)-\mathrm{conf}(B_m)\bigr|.
\end{equation}
\emph{Temperature scaling} (Guo et al., 2017) is the cheapest fix. It divides the logits
by a single scalar $T$ fitted on the validation set to minimise the negative
log-likelihood. Because it does not change the $\argmax$, accuracy is untouched. This is
the same temperature that controls sampling diversity in \cref{ch:pretrain,ch:scaling},
used here for the opposite purpose of making probabilities honest.

\begin{lstlisting}[caption={Fitting a calibration temperature on cached validation logits.},label={lst:classify:temp}]
def fit_temperature(val_logits, val_labels, steps=200, lr=0.01):
    # val_logits: (n, C), collected once under torch.no_grad()
    log_t = torch.zeros(1, requires_grad=True)          # T = exp(log_t) > 0
    opt = torch.optim.Adam([log_t], lr=lr)
    for _ in range(steps):
        loss = F.cross_entropy(val_logits / log_t.exp(), val_labels)
        opt.zero_grad(); loss.backward(); opt.step()
    return log_t.exp().item()
\end{lstlisting}

Note the systems shape of this step. We run the expensive model \emph{once} over the
validation set, cache an $n\times C$ matrix of logits, and fit on that. The fit
takes milliseconds. The same pattern of caching model outputs and tuning cheap
post-processing on top appears again with verifiers and scorers in
\cref{ch:eval,ch:refine}.

\subsection{Cascades: spend big-model compute only on hard inputs}

Suppose a large instruction-tuned model (\cref{ch:instruct}) classifies more accurately
than our fine-tuned GPT-2, but costs $R$ times more per token. A \emph{cascade} runs the
small classifier on everything and forwards a request to the large model only when the
small model's calibrated confidence $\max_c p(c\mid x)$ is below a threshold. If a
fraction $e$ of traffic is escalated, the expected cost per request in units of the
small model is
\begin{equation}
\text{cost}(e) = 1 + e\,R,\qquad \frac{\text{cost}(e)}{\text{cost of large only}} = \frac{1+eR}{R}.
\label{eq:classify:cascade}
\end{equation}
Calibration matters here: with an over-confident small model, the escalated set misses
exactly the inputs it gets wrong.

\begin{figure}[ht]
\centering
\begin{minipage}{0.44\linewidth}\centering
\begin{tikzpicture}[box/.style={draw=#1,fill=#1!12,rounded corners=2pt,minimum height=0.8cm,
    align=center,font=\scriptsize,text width=1.9cm},arr/.style={-{Latex[length=1.8mm]},thick}]
\node[box=srcgray] (in) at (0,0) {request};
\node[box=keygreen] (small) at (0,-1.3) {small classifier\\(GPT-2 FT)};
\node[draw=sysblue,diamond,aspect=1.8,inner sep=1pt,font=\scriptsize,align=center] (dec) at (0,-2.75)
  {$\max_c p \ge \tau$?};
\node[box=sysblue] (ans) at (-1.4,-4.3) {answer\\(fraction $1-e$)};
\node[box=warnred] (big) at (1.4,-4.3) {large LLM\\(fraction $e$)};
\draw[arr] (in)--(small); \draw[arr] (small)--(dec);
\draw[arr] (dec.west) -| node[pos=0.25,above,font=\scriptsize] {yes} (ans.north);
\draw[arr] (dec.east) -| node[pos=0.25,above,font=\scriptsize] {no} (big.north);
\end{tikzpicture}
\end{minipage}\hfill
\begin{minipage}{0.54\linewidth}\centering
\begin{tikzpicture}
\begin{axis}[width=\linewidth,height=5cm,xlabel={escalated fraction $e$},
  ylabel={cost / large-only},ymin=0,ymax=1.05,xmin=0,xmax=1,grid=major,grid style={gray!20},
  legend style={font=\scriptsize,at={(0.98,0.03)},anchor=south east}]
\addplot[thick,warnred,domain=0:1,samples=50] {(1+x*94)/94};
\addlegendentry{$R=94$ (8B vs.\ GPT-2)}
\addplot[thick,sysblue,dashed,domain=0:1,samples=50] {(1+x*10)/10};
\addlegendentry{$R=10$}
\end{axis}
\end{tikzpicture}
\end{minipage}
\caption{Left: a two-stage classifier cascade. Right: the relative cost from
\cref{eq:classify:cascade}. When the large model is $R\approx94\times$ more expensive per
token (an 8B-parameter model versus GPT-2 small's 85M non-embedding parameters,
prefill only), escalating 5\% of traffic costs about 6\% of the large-only budget.}
\label{fig:classify:cascade}
\end{figure}

\begin{napkinmath}[Pricing a cascade]
Prefill FLOPs per token scale with non-embedding parameters: $2\times 85\text{M}$ for
GPT-2 small versus about $2\times 8\text{B}$ for an 8B model, so $R\approx94$. If the
small model is confident on 95\% of messages, $\text{cost}=1+0.05\times94=5.7$ small-model
units, against 94 for the large model alone: a $16\times$ saving. And if the large
model must \emph{generate} a verdict instead of scoring label tokens in a single
prefill, $R$ grows further, because decoding is memory-bound (\cref{ch:inference}).
\end{napkinmath}

\subsection{One trunk, many heads}

Platforms often run several classifiers over every message: spam, abuse, language,
intent, PII. If they all fine-tune only the top of the same base model, the frozen
bottom of the network computes \emph{identical} activations for each of them. Sharing
that computation is pure profit.

\begin{figure}[ht]
\centering
\begin{tikzpicture}[blk/.style={draw=#1,fill=#1!12,rounded corners=1.5pt,minimum height=0.55cm,
    font=\scriptsize,align=center},arr/.style={-{Latex[length=1.6mm]},srcgray}]
\node[blk=srcgray,minimum width=3.4cm] (trunk) at (0,0) {shared frozen blocks 1--11\\(run once per message)};
\foreach \k/\name/\c in {0/spam/keygreen,1/abuse/warnred,2/intent/sysblue,3/language/labpurple}{
  \node[blk=\c,minimum width=2.2cm] (b\k) at (4.4,1.5-\k) {block 12 (\name)};
  \node[blk=\c,minimum width=1.9cm] (h\k) at (7.1,1.5-\k) {head (\name)};
  \draw[arr] (trunk.east) -- (b\k.west); \draw[arr] (b\k) -- (h\k);}
\end{tikzpicture}
\caption{Serving several fine-tuned classifiers that share a frozen trunk. Each task
owns only its top block and head (about \SI{14}{MB} in bf16 for GPT-2 small).}
\label{fig:classify:multihead}
\end{figure}

With $K$ tasks that each fine-tune only the last of $L$ blocks, separate models evaluate
$KL$ blocks per token, while a shared trunk evaluates $(L-1)+K$. For $L=12$ and $K=8$
that is $96$ versus $19$, a $5\times$ saving in FLOPs \emph{and} in weight memory traffic.
If the tasks only train heads (\code{trainable\_blocks=0}), the saving approaches
$K\times$. LoRA gives the same economics for adapters spread through the whole network,
using the batched multi-adapter kernels of \cref{sec:instruct:multilora}.

\subsection{Classifying with a modern decoder}

Nothing in this chapter is specific to GPT-2. The Qwen3 model of \cref{ch:modern} uses
the same attribute names (\code{trf\_blocks}, \code{final\_norm}, \code{out\_head}), so
\code{add\_classification\_head} applies unchanged. Three differences deserve care:
\begin{itemize}
\item \textbf{Tied embeddings.} In Qwen3-0.6B, \code{out\_head.weight} \emph{is}
\code{tok\_emb.weight}. Replacing \code{out\_head} breaks the tie, which is what we want.
The input embedding stays frozen and the new head is independent. But parameter counts
and checkpoint sizes change, and code that re-ties weights on load must be skipped.
\item \textbf{A real padding token.} Qwen's tokenizer has distinct
\code{<|endoftext|>} and \code{<|im\_end|>} tokens, so padding never collides with an
end-of-turn marker that the model treats as meaningful.
\item \textbf{RoPE makes left padding exact.} Attention scores depend only on relative
positions, so left padding plus a key mask gives the same last-token logits as an
unpadded forward pass, up to floating-point summation order. The core test suite checks
this property for batched generation.
\end{itemize}

\begin{lstlisting}[caption={The same recipe on a Qwen3-shaped model (tiny config shown).},label={lst:classify:qwen}]
from llmsys.qwen3 import Qwen3Model, QWEN3_CONFIG_TINY
from llmsys.finetune import add_classification_head, pad_batch

model = add_classification_head(Qwen3Model(QWEN3_CONFIG_TINY), num_classes=2)
ids, lengths = pad_batch([[5, 6, 7, 8], [9, 10]], pad_id=0, side="left")
pad_mask = ids != 0                              # True for real tokens
logits = model(ids, pad_mask=pad_mask)[:, -1]   # (2, 2): last position is real
\end{lstlisting}

% =====================================================================================
\section{Summary}\label{sec:classify:summary}
\begin{itemize}
\item Classification fine-tuning replaces the vocabulary projection with a $C$-way head
and trains on labelled examples. It is cheaper and easier to evaluate than instruction
fine-tuning, but the resulting model does only that one task.
\item Data preparation matters: balance or reweight classes, deduplicate before
splitting, and split with a fixed seed into train/validation/test (70/10/20 here).
\item Variable-length inputs must be padded. Padding to the global maximum wastes most
of the compute; dynamic padding plus length bucketing removes nearly all of it.
\item Under a causal mask only the last position has seen the whole input, so the head
reads that position. With right padding, gather the last \emph{real} token, or use left
padding with a key mask.
\item Freezing the trunk removes optimizer state and, when the frozen layers are at the
bottom, also their backward pass and stored activations. It even allows caching trunk
features across epochs.
\item Serving a classifier is pure prefill: compute-bound, no decode loop, no vocabulary
head. Padding-aware batching is the main lever on throughput.
\item Thresholds should be chosen on calibrated probabilities. Calibrated small
classifiers enable cascades and shared-trunk multi-task serving, which cut cost by an
order of magnitude.
\end{itemize}

% =====================================================================================
\section{Exercises}\label{sec:classify:exercises}

\begin{exercise}[Global versus dynamic padding]
Train the classifier twice: once with every split padded to the longest training
message (120 tokens), and once with per-batch dynamic padding and last-real-token
pooling. Compare accuracy and wall-clock time per epoch. Then pad everything to 1{,}024
tokens and explain what happens to time and to accuracy when the head reads position
$-1$.
\end{exercise}

\begin{exercise}[How many layers to unfreeze]\label{ex:classify:layers}
Sweep \code{trainable\_blocks}$\in\{0,1,2,4,12\}$ (0 = head and final norm only). For
each setting, report test accuracy, trainable parameter count, peak memory, and
time per step. Predict the time per step from the $2N + 4N_{\text{tr}}$ model before
measuring, and explain any discrepancy.
\end{exercise}

\begin{exercise}[Trunk caching]
With \code{trainable\_blocks=1}, precompute the outputs of blocks 1--11 for the whole
training set and train block 12 and the head on the cached tensors. Verify that the
loss trajectory matches uncached training (dropout off). How large is the cache, and
how much faster is each epoch?
\end{exercise}

\begin{exercise}[Class weights instead of undersampling]
Use all 5{,}574 messages with the original 87/13 class ratio. Compare (a) plain
cross-entropy, (b) cross-entropy with class weights inversely proportional to class
frequency, and (c) undersampling. Report accuracy, precision, recall and F1 for the
spam class. Which metric would you put on a dashboard, and why?
\end{exercise}

\begin{exercise}[Padding invariance]
Prove that with right padding, causal attention, and last-real-token pooling, a
message's logits do not depend on how much padding follows it. Then show with a
counterexample that reading position $-1$ does depend on padding amount. (Hint: what
does the padding position attend to?)
\end{exercise}

\begin{exercise}[Left padding and absolute positions]
Classify a batch with left padding and a key-padding mask using GPT-2's learned
absolute position embeddings. Measure how much the logits differ from the unpadded
single-message logits, and explain why. Would the result change for a RoPE model?
\end{exercise}

\begin{exercise}[Token-budget batching]
Implement a sampler that yields batches with at most $K$ tokens after padding (e.g.\
$K=2048$). Compare the number of batches, padding waste, and step-time variance against
fixed-size bucketed batches.
\end{exercise}

\begin{exercise}[Zero-shot baseline]
Without any fine-tuning, classify messages by comparing $\log p(\text{`` spam''})$ and
$\log p(\text{`` ham''})$ after a prompt such as
``Message: \dots\textbackslash nCategory:''. Report accuracy on the test split. How
does the gap to the fine-tuned model change with model size (GPT-2 small, medium,
large)?
\end{exercise}

\begin{exercise}[Serving capacity]
A service must classify 20{,}000 messages per second (mean 35 tokens, p99 160 tokens)
with p99 latency under \SI{50}{ms}. Using the napkin model of
\cref{sec:classify:use}, estimate how many L4 GPUs are needed with (a) random batching
and (b) bucketing. Which assumption in your estimate is most likely to be wrong?
\end{exercise}

\begin{exercise}[Deduplication]
Count exact duplicates in the SMS collection, and near-duplicates after lowercasing and
replacing digit runs with a placeholder. Re-split after deduplication and compare test
accuracy to the original split. What does the difference tell you?
\end{exercise}

\begin{exercise}[Calibration]
Collect validation logits once, fit a temperature with \cref{lst:classify:temp}, and
report ECE (15 bins) on the test split before and after. Show analytically that
temperature scaling cannot change accuracy.
\end{exercise}

\begin{exercise}[Cascade design]
Using your calibrated classifier, plot escalation rate $e(\tau)$ and cascade accuracy
against $\tau$, assuming the large model is always right. Find the cheapest $\tau$
that reaches 99.5\% accuracy and compute its cost with \cref{eq:classify:cascade}.
\end{exercise}

\begin{exercise}[Layer-wise probes]\label{ex:classify:probe}
Use \cref{lst:classify:probe} to collect last-real-token features after every block of
pretrained GPT-2 small on the balanced SMS splits, fit a logistic-regression probe per
layer, and plot validation accuracy against depth. At which depth would you cut the
network for serving, and what is the predicted FLOP saving?
\end{exercise}

% =====================================================================================
\section{Lab 16: Classification fine-tuning and padding-aware batching}\label{sec:classify:lab}

\begin{labbox}{16}{Classification fine-tuning and padding-aware batching}
\labdir{lab16_classify}
\labsection{Goals}
Build the full classification fine-tuning pipeline: balancing, splitting, dynamic
padding, the head swap with selective freezing, padding-invariant pooling, and the
training loop. Then measure how much compute padding wastes and how much bucketing
recovers.

\labsection{Setup}
\code{cd code; LAB\_IMPL=starter python -m pytest -q labs/lab09\_classify}.
The tests run on CPU in under a minute using a tiny GPT and a synthetic spam task
(\code{make\_synthetic\_spam}: ``spam'' messages contain marker tokens, and lengths are
log-normal). \code{run.py} repeats the experiment with GPT-2 small on the real SMS Spam
Collection on your own machine.

\labsection{Tasks}
\begin{enumerate}
\item \code{balance\_by\_undersampling}, \code{split\_indices}: seeded, exact sizes,
every index exactly once.
\item \code{pad\_batch}: left/right padding, optional truncation, returns lengths.
\item \code{add\_classification\_head}: freeze everything, swap the head, unfreeze the
last $k$ blocks and the final norm.
\item \code{select\_class\_logits} and \code{accuracy\_loader}: last-real-token
pooling. The test checks that predictions do not change when more padding follows a
message.
\item \code{length\_bucketed\_batches} and \code{padding\_stats}: sort-within-mega-batch
bucketing and a waste meter. The test requires bucketing to cut waste by at least
$3\times$.
\item \code{train\_classifier}: AdamW over trainable parameters, training accuracy
accumulated for free, validation accuracy per epoch. The test requires $>80\%$ test
accuracy on the synthetic task.
\end{enumerate}

\labsection{Measurements}
\begin{itemize}
\item Reproduce \cref{fig:classify:waste} for your own length distribution (synthetic
and, if you have it, the real SMS data with the GPT-2 tokenizer).
\item With \code{run.py} on a GPU: time per epoch and peak memory for
\code{trainable\_blocks}$\in\{1,12\}$, with and without \code{--bucketed}. Expect
bucketing to cut epoch time roughly in proportion to the waste reduction once batches
are large enough to be compute-bound. Expect \code{trainable\_blocks=1} to be
roughly $2$--$3\times$ faster than full fine-tuning per step.
\item Accuracy before and after fine-tuning: near 50\% before, mid-to-high 90s after on
real SMS data.
\end{itemize}

\labsection{Deliverables}
Passing tests; a short table (accuracy, time/epoch, peak memory, trainable params) for
the configurations above; one paragraph explaining the measured speed-up from freezing
in terms of forward and backward FLOPs.

\labsection{Stretch goals}
\begin{itemize}
\item Implement trunk caching and token-budget batching (Exercises 16.3 and 9.7).
\item Serve the classifier behind a minimal HTTP endpoint with a micro-batching queue:
collect requests for up to \SI{5}{ms} or 64 messages, bucket by length, and run one
forward pass. Plot throughput against p99 latency as the batching window varies.
\end{itemize}
\end{labbox}
